\textbf{1.\ (Binary stones).}
Consider the following game: two players may take from a heap any number of stones that is a power of \(2\) (including \(1 = 2^0\)).
The player who has taken the last stone wins.
For which number of stones does the first player win?
How should he play?
(You should not only specify the strategy but also prove that it works.)

\textbf{Solution.}

Let \(n\) be the number of stones.
Allowed moves are to remove \(1, 2, 4, 8, \dots, 2^k\) stones.
We seek the losing positions:

\medskip
\noindent\textbf{Small cases.}
\[
\begin{array}{c|cccccccccc}
n & 0 & 1 & 2 & 3 & 4 & 5 & 6 & 7 & 8 & 9\\\hline
\text{outcome} & \text{lose} & \text{win} & \text{win} & \text{lose} & \text{win} & \text{win} & \text{lose} & \text{win} & \text{win} & \text{lose}
\end{array}
\]
The losing positions among these are \(0, 3, 6, 9, \dots\)

\medskip
\noindent\textbf{Claim.}
The losing positions are exactly the multiples of \(3\).
Equivalently, the first player wins if and only if \(n \not\equiv 0 \pmod{3}\). For example, 1, 2 or 4 and etc.

\medskip
\noindent\textbf{Proof.}
Powers of \(2\) modulo \(3\) alternate:
\[
2^0 \equiv 1\pmod{3} \qquad 2^1 \equiv 2\pmod{3}
\]
\[
2^2 \equiv 1 \pmod{3} \qquad 2^3 \equiv 2 \pmod{3}
\]
In general \(2^k \equiv (-1)^k \pmod{3}\), so every allowed move removes either \(1\) or \(2\) modulo \(3\), never \(0\) modulo \(3\).

\smallskip
If \(n \equiv 0 \pmod{3}\), then for any move \(n \mapsto n - 2^k\) we have \(n - 2^k \not\equiv 0 \pmod{3}\).

If \(n \not\equiv 0 \pmod{3}\), there is a move to a multiple of \(3\):

\begin{center}

\noindent\textbf{1.} \(n \equiv 1 \pmod{3}\), remove \(1 = 2^0\)

\noindent\textbf{2.} \(n \equiv 2 \pmod{3}\), remove \(2 = 2^1\)

\end{center}

Thus from a multiple of \(3\) every move goes to a non-multiple of \(3\), and from a non-multiple of \(3\) there is a move to a multiple of \(3\).
The terminal position \(n = 0\) is a multiple of \(3\) and is losing for the player to move.
Hence multiples of \(3\). Therefore the first player wins if and only if \(n \not\equiv 0 \pmod{3}\).

\medskip
\noindent\textbf{Optimal strategy.}
If \(n \equiv 1 \pmod{3}\), remove \(1\) stone; if \(n \equiv 2 \pmod{3}\), remove \(2\) stones.
Thereafter always leave a heap size divisible by \(3\).
Since the opponent can only subtract \(1\) or \(2\) modulo \(3\) from the heap, you can always restore divisibility by \(3\) on your turn.
Eventually the opponent faces \(n = 0\) and loses.
